%% lecture_outline.tex -- the ONE place the eleven lecture titles and their
%% subchapters are written down.
%%
%% Why this file exists. The orientation deck carries a "lecture in one slide"
%% for the whole course, and every chapter now carries the same card for its
%% own lecture. Two copies of the same list drift: ch01's first subchapter was
%% renamed to "The small part and the big picture" while ch00 still announced
%% "When to Use Machine Learning". So there are no copies. Both render from
%% here, and tools/check_lecture_outline.sh fails the build if this file and
%% the chapters' own \subsection commands disagree.
%%
%% To change a subchapter: rename the \subsection in the chapter file AND the
%% matching line here. The checker tells you when you did only one of them.
%% "Summary" subchapters are deliberately left out of the cards; the checker
%% ignores them too.

\makeatletter

% \outlinetitle{<roman>} -- the lecture's title, the book's own chapter title.
\newcommand{\outlinetitle}[1]{\csname td@lt@#1\endcsname}
% \outlineitems{<roman>} -- its subchapters, already formatted as table rows.
\newcommand{\outlineitems}[1]{\csname td@lo@#1\endcsname}

\newcommand{\td@setlecture}[3]{%
  \expandafter\gdef\csname td@lt@#1\endcsname{#2}%
  \expandafter\gdef\csname td@lo@#1\endcsname{#3}%
}

\td@setlecture{I}{Overview of Machine Learning Systems}{%
  The small part and the big picture\\[0.2em]
  Understanding Machine Learning Systems\\[0.2em]
}

\td@setlecture{II}{Introduction to Machine Learning Systems Design}{%
  Business and ML Objectives\\[0.2em]
  Requirements for ML Systems\\[0.2em]
  Iterative Process\\[0.2em]
  Framing ML Problems\\[0.2em]
  Mind Versus Data\\[0.2em]
}

\td@setlecture{III}{Data Engineering Fundamentals}{%
  Data Sources\\[0.2em]
  Data Formats\\[0.2em]
  Data Models\\[0.2em]
  Data Storage Engines and Processing\\[0.2em]
  Modes of Dataflow\\[0.2em]
  Batch Processing Versus Stream Processing\\[0.2em]
}

\td@setlecture{IV}{Training Data}{%
  Sampling\\[0.2em]
  Labeling\\[0.2em]
  Class Imbalance\\[0.2em]
  Data Augmentation\\[0.2em]
}

\td@setlecture{V}{Feature Engineering}{%
  Learned Features Versus Engineered Features\\[0.2em]
  Common Feature Engineering Operations\\[0.2em]
  Data Leakage\\[0.2em]
  Engineering Good Features\\[0.2em]
}

\td@setlecture{VI}{Model Development and Offline Evaluation}{%
  Model Development and Training\\[0.2em]
  Model Offline Evaluation\\[0.2em]
}

\td@setlecture{VII}{Model Deployment and Prediction Service}{%
  Machine Learning Deployment Myths\\[0.2em]
  Batch Prediction Versus Online Prediction\\[0.2em]
  Model Compression\\[0.2em]
  ML on the Cloud and on the Edge\\[0.2em]
}

\td@setlecture{VIII}{Data Distribution Shifts and Monitoring}{%
  Causes of ML System Failures\\[0.2em]
  Data Distribution Shifts\\[0.2em]
  Monitoring and Observability\\[0.2em]
}

\td@setlecture{IX}{Continual Learning and Test in Production}{%
  Continual Learning\\[0.2em]
  Test in Production\\[0.2em]
}

\td@setlecture{X}{Infrastructure and Tooling for MLOps}{%
  Storage and Compute\\[0.2em]
  Development Environment\\[0.2em]
  Resource Management\\[0.2em]
  ML Platform\\[0.2em]
  Build Versus Buy\\[0.2em]
}

\td@setlecture{XI}{The Human Side of Machine Learning}{%
  User Experience\\[0.2em]
  Team Structure\\[0.2em]
  Responsible AI\\[0.2em]
}

% \outlinecard{<roman>} -- one lecture as a booktabs card: numeral and title in
% the head, subchapters under it. Three of them side by side fill the
% orientation frames; one on its own opens a chapter.
\newcommand{\outlinecard}[1]{%
  \footnotesize
  \renewcommand{\arraystretch}{1.1}%
  \begin{tabular}[t]{@{}>{\raggedright\arraybackslash}p{\linewidth}@{}}
    \toprule
    \parbox[t][2\baselineskip][t]{\linewidth}{\raggedright
      \alertbf{#1}\quad \textbf{\outlinetitle{#1}}}\\
    \midrule
    \outlineitems{#1}
    \bottomrule
  \end{tabular}%
}

\makeatother
